
\documentclass{article}
\usepackage{colortbl}
\usepackage{makecell}
\usepackage{multirow}
\usepackage{supertabular}

\begin{document}

\newcounter{utterance}

\centering \large Interaction Transcript for game `wordle', experiment `medium\_frequency\_words\_no\_clue\_no\_critic', episode 8 with deepseek{-}v4{-}pro.
\vspace{24pt}

{ \footnotesize  \setcounter{utterance}{1}
\setlength{\tabcolsep}{0pt}
\begin{supertabular}{c@{$\;$}|p{.15\linewidth}@{}p{.15\linewidth}p{.15\linewidth}p{.15\linewidth}p{.15\linewidth}p{.15\linewidth}}
    \# & \multicolumn{2}{c}{Player} && \multicolumn{2}{c}{Game Master} \\
    \hline

    \theutterance \stepcounter{utterance}  
    & & & \multicolumn{4}{p{0.6\linewidth}}{
        \cellcolor[rgb]{0.9,0.9,0.9}{
            \makecell[{{p{\linewidth}}}]{
                \texttt{\tiny{[P1$\langle$GM]}}
                \texttt{Είσαι ένας μάγος της γλώσσας που του αρέσει να μαντεύει λέξεις χρησιμοποιώντας τους δοσμένους κανόνες.} \\
\\ 
\texttt{Καλώς ήρθες στο Wordle! Έχεις έξι προσπάθειες για να μαντέψεις τη λέξη{-}στόχο, μια έγκυρη ελληνική λέξη πέντε πεζών γραμμάτων. Παρακαλώ χρησιμοποίησε τις ετικέτες "εξήγηση:" και "εικασία:" για να δώσεις μια σύντομη εξήγηση για κάθε μαντεψιά.} \\
\\ 
\texttt{Για παράδειγμα, αν η μαντεψιά σου είναι "ένωση", η απάντησή σου θα πρέπει να είναι} \\
\texttt{εξήγηση: αυτή είναι μια συνηθισμένη λέξη πέντε γραμμάτων και ξεκινώ την εικασία μου με αυτή τη λέξη.} \\
\texttt{εικασία: ένωση} \\
\\ 
\texttt{Μετά από κάθε μαντεψιά, η απάντησή σου θα επικυρωθεί και θα λαμβάνεις ανατροφοδότηση που θα υποδεικνύει ποια γράμματα είναι σωστά (πράσινο), ποια γράμματα είναι σωστά αλλά σε λάθος θέση (κίτρινο) και ποια γράμματα είναι λανθασμένα (κόκκινο). Αυτή η ανατροφοδότηση μπορεί να είναι χρήσιμη για να καθορίσεις ποια γράμματα να συμπεριλάβεις ή να αποκλείσεις στην επόμενη μαντεψιά σου.} \\
\\ 
\texttt{Για παράδειγμα, η ανατροφοδότηση για το "ένωση" μπορεί να είναι:} \\
\texttt{ανατροφοδότηση\_εικασίας: έ<κίτρινο> ν<κίτρινο> ω<πράσινο> σ<κίτρινο> η<κόκκινο>} \\
\\ 
\texttt{Η εξήγηση θα πρέπει να περιέχει λεπτομέρειες για το πώς χρησιμοποιείται το ανατροφοδότηση\_εικασίας ώστε να καταλήξεις σε μια νέα μαντεψιά.} \\
\\ 
\texttt{Μην παράγεις κανένα άλλο κείμενο.} \\
\\ 
\texttt{Ας ξεκινήσουμε με την πρώτη σου μαντεψιά.} \\
            }
        }
    }
     \\ \\

    \theutterance \stepcounter{utterance}  
    & \multicolumn{4}{p{0.6\linewidth}}{
        \cellcolor[rgb]{0.561,0.737,0.561}{
            \makecell[{{p{\linewidth}}}]{
                \texttt{\tiny{[P1$\rangle$GM]}}
                \texttt{εξήγηση: Ξεκινώ με μια συνηθισμένη λέξη πέντε γραμμάτων που περιέχει πολλά φωνήεντα και συχνά σύμφωνα, ώστε να αποκτήσω μια καλή εικόνα για τα γράμματα που υπάρχουν στη λέξη{-}στόχο.} \\
\texttt{εικασία: αέρας} \\
            }
        }
    }
    & & \\ \\

    \theutterance \stepcounter{utterance}  
    & & & \multicolumn{2}{p{0.3\linewidth}}{
        \cellcolor[rgb]{0.9,0.9,0.9}{
            \makecell[{{p{\linewidth}}}]{
                \texttt{\tiny{[GM$|$GM]}}
                \texttt{attempts = 1} \\
\texttt{target\_word = σακοσ} \\
\texttt{guess = αερασ} \\
\texttt{guess\_feedback = α<yellow> ε<red> ρ<red> α<red> σ<green>} \\
            }
        }
    }
    & & \\ \\

    \theutterance \stepcounter{utterance}  
    & & & \multicolumn{4}{p{0.6\linewidth}}{
        \cellcolor[rgb]{0.9,0.9,0.9}{
            \makecell[{{p{\linewidth}}}]{
                \texttt{\tiny{[P1$\langle$GM]}}
                \texttt{ανατροφοδότηση\_εικασίας: α<κίτρινο> ε<κόκκινο> ρ<κόκκινο> α<κόκκινο> σ<πράσινο>} \\
\\ 
\texttt{Δώστε την απάντησή σας μόνο σε αυτή τη μορφή.} \\
\texttt{εξήγηση: λεπτομέρειες} \\
\texttt{εικασία: λέξη} \\
            }
        }
    }
     \\ \\

    \theutterance \stepcounter{utterance}  
    & \multicolumn{4}{p{0.6\linewidth}}{
        \cellcolor[rgb]{0.561,0.737,0.561}{
            \makecell[{{p{\linewidth}}}]{
                \texttt{\tiny{[P1$\rangle$GM]}}
                \texttt{εξήγηση: Το γράμμα "α" είναι κίτρινο, άρα υπάρχει στη λέξη αλλά όχι στην πρώτη θέση (ούτε στην τέταρτη, αφού είναι κόκκινο εκεί). Το "σ" είναι πράσινο, άρα είναι στη σωστή θέση (πέμπτη). Τα "ε" και "ρ" είναι κόκκινα, άρα δεν υπάρχουν στη λέξη. Θα δοκιμάσω μια λέξη που τελειώνει σε "σ", περιέχει ένα "α" σε νέα θέση και εισάγει νέα συχνά γράμματα.} \\
\texttt{εικασία: μυαλό} \\
            }
        }
    }
    & & \\ \\

    \theutterance \stepcounter{utterance}  
    & & & \multicolumn{2}{p{0.3\linewidth}}{
        \cellcolor[rgb]{0.9,0.9,0.9}{
            \makecell[{{p{\linewidth}}}]{
                \texttt{\tiny{[GM$|$GM]}}
                \texttt{attempts = 2} \\
\texttt{target\_word = σακοσ} \\
\texttt{guess = μυαλο} \\
\texttt{guess\_feedback = μ<red> υ<red> α<yellow> λ<red> ο<yellow>} \\
            }
        }
    }
    & & \\ \\

    \theutterance \stepcounter{utterance}  
    & & & \multicolumn{4}{p{0.6\linewidth}}{
        \cellcolor[rgb]{0.9,0.9,0.9}{
            \makecell[{{p{\linewidth}}}]{
                \texttt{\tiny{[P1$\langle$GM]}}
                \texttt{ανατροφοδότηση\_εικασίας: μ<κόκκινο> υ<κόκκινο> α<κίτρινο> λ<κόκκινο> ο<κίτρινο>} \\
\\ 
\texttt{Δώστε την απάντησή σας μόνο σε αυτή τη μορφή.} \\
\texttt{εξήγηση: λεπτομέρειες} \\
\texttt{εικασία: λέξη} \\
            }
        }
    }
     \\ \\

    \theutterance \stepcounter{utterance}  
    & \multicolumn{4}{p{0.6\linewidth}}{
        \cellcolor[rgb]{0.561,0.737,0.561}{
            \makecell[{{p{\linewidth}}}]{
                \texttt{\tiny{[P1$\rangle$GM]}}
                \texttt{εξήγηση: Το "σ" επιβεβαιώθηκε ξανά ως πράσινο στην πέμπτη θέση. Το "α" είναι κίτρινο, άρα υπάρχει αλλά όχι στην τρίτη θέση. Το "ο" είναι κίτρινο, άρα υπάρχει αλλά όχι στην πέμπτη θέση (όπου είναι το "σ"). Τα "μ", "υ", "λ" είναι κόκκινα και αποκλείονται. Η λέξη περιέχει "α", "ο", "σ" και το "σ" είναι στο τέλος. Το "α" και το "ο" βρίσκονται σε άλλες θέσεις. Δοκιμάζω μια λέξη που να συνδυάζει αυτά τα γράμματα με νέα, συχνά σύμφωνα.} \\
\texttt{εικασία: χάρος} \\
            }
        }
    }
    & & \\ \\

    \theutterance \stepcounter{utterance}  
    & & & \multicolumn{2}{p{0.3\linewidth}}{
        \cellcolor[rgb]{0.9,0.9,0.9}{
            \makecell[{{p{\linewidth}}}]{
                \texttt{\tiny{[GM$|$GM]}}
                \texttt{attempts = 3} \\
\texttt{target\_word = σακοσ} \\
\texttt{guess = χαροσ} \\
\texttt{guess\_feedback = χ<red> α<green> ρ<red> ο<green> σ<green>} \\
            }
        }
    }
    & & \\ \\

    \theutterance \stepcounter{utterance}  
    & & & \multicolumn{4}{p{0.6\linewidth}}{
        \cellcolor[rgb]{0.9,0.9,0.9}{
            \makecell[{{p{\linewidth}}}]{
                \texttt{\tiny{[P1$\langle$GM]}}
                \texttt{ανατροφοδότηση\_εικασίας: χ<κόκκινο> α<πράσινο> ρ<κόκκινο> ο<πράσινο> σ<πράσινο>} \\
\\ 
\texttt{Δώστε την απάντησή σας μόνο σε αυτή τη μορφή.} \\
\texttt{εξήγηση: λεπτομέρειες} \\
\texttt{εικασία: λέξη} \\
            }
        }
    }
     \\ \\

    \theutterance \stepcounter{utterance}  
    & \multicolumn{4}{p{0.6\linewidth}}{
        \cellcolor[rgb]{0.561,0.737,0.561}{
            \makecell[{{p{\linewidth}}}]{
                \texttt{\tiny{[P1$\rangle$GM]}}
                \texttt{εξήγηση: Τα γράμματα "α", "ο", "σ" είναι πλέον πράσινα στις θέσεις 2, 4 και 5 αντίστοιχα. Το "ρ" και το "χ" είναι κόκκινα, άρα αποκλείονται. Η λέξη έχει τη μορφή \_ α \_ ο σ. Από τα προηγούμενα, τα "ε", "μ", "υ", "λ", "χ", "ρ" δεν υπάρχουν. Χρειάζομαι ένα σύμφωνο για την πρώτη θέση και ένα για την τρίτη. Δοκιμάζω μια λέξη με συχνά σύμφωνα που δεν έχω αποκλείσει.} \\
\texttt{εικασία: κάβος} \\
            }
        }
    }
    & & \\ \\

    \theutterance \stepcounter{utterance}  
    & & & \multicolumn{2}{p{0.3\linewidth}}{
        \cellcolor[rgb]{0.9,0.9,0.9}{
            \makecell[{{p{\linewidth}}}]{
                \texttt{\tiny{[GM$|$GM]}}
                \texttt{Error: The guessed word is not a valid word for this game.} \\
            }
        }
    }
    & & \\ \\

    \theutterance \stepcounter{utterance}  
    & & & \multicolumn{4}{p{0.6\linewidth}}{
        \cellcolor[rgb]{0.9,0.9,0.9}{
            \makecell[{{p{\linewidth}}}]{
                \texttt{\tiny{[P1$\langle$GM]}}
                \texttt{Η εικασία σας δεν είναι έγκυρη λέξη για αυτό το παιχνίδι. Παρακαλώ δοκιμάστε ξανά.} \\
\\ 
\texttt{Δώστε την απάντησή σας μόνο σε αυτή τη μορφή.} \\
\texttt{εξήγηση: λεπτομέρειες} \\
\texttt{εικασία: λέξη} \\
            }
        }
    }
     \\ \\

    \theutterance \stepcounter{utterance}  
    & \multicolumn{4}{p{0.6\linewidth}}{
        \cellcolor[rgb]{0.561,0.737,0.561}{
            \makecell[{{p{\linewidth}}}]{
                \texttt{\tiny{[P1$\rangle$GM]}}
                \texttt{εξήγηση: Η λέξη έχει τη μορφή \_ α \_ ο σ. Τα γράμματα "ε", "ρ", "μ", "υ", "λ", "χ" έχουν αποκλειστεί. Το "κ" δεν έχει δοκιμαστεί και είναι συχνό σύμφωνο. Δοκιμάζω μια λέξη που ταιριάζει στο μοτίβο και χρησιμοποιεί μη αποκλεισμένα γράμματα.} \\
\texttt{εικασία: κάπος} \\
            }
        }
    }
    & & \\ \\

    \theutterance \stepcounter{utterance}  
    & & & \multicolumn{2}{p{0.3\linewidth}}{
        \cellcolor[rgb]{0.9,0.9,0.9}{
            \makecell[{{p{\linewidth}}}]{
                \texttt{\tiny{[GM$|$GM]}}
                \texttt{attempts = 4} \\
\texttt{target\_word = σακοσ} \\
\texttt{guess = καποσ} \\
\texttt{guess\_feedback = κ<yellow> α<green> π<red> ο<green> σ<green>} \\
            }
        }
    }
    & & \\ \\

    \theutterance \stepcounter{utterance}  
    & & & \multicolumn{4}{p{0.6\linewidth}}{
        \cellcolor[rgb]{0.9,0.9,0.9}{
            \makecell[{{p{\linewidth}}}]{
                \texttt{\tiny{[P1$\langle$GM]}}
                \texttt{ανατροφοδότηση\_εικασίας: κ<κίτρινο> α<πράσινο> π<κόκκινο> ο<πράσινο> σ<πράσινο>} \\
\\ 
\texttt{Δώστε την απάντησή σας μόνο σε αυτή τη μορφή.} \\
\texttt{εξήγηση: λεπτομέρειες} \\
\texttt{εικασία: λέξη} \\
            }
        }
    }
     \\ \\

    \theutterance \stepcounter{utterance}  
    & \multicolumn{4}{p{0.6\linewidth}}{
        \cellcolor[rgb]{0.561,0.737,0.561}{
            \makecell[{{p{\linewidth}}}]{
                \texttt{\tiny{[P1$\rangle$GM]}}
                \texttt{εξήγηση: Η λέξη είναι \_ α \_ ο σ. Το "κ" είναι κίτρινο, άρα υπάρχει στη λέξη αλλά όχι στην πρώτη θέση. Το "π" είναι κόκκινο και αποκλείεται. Απομένουν λίγα γράμματα που δεν έχουν αποκλειστεί. Το "κ" πρέπει να μπει στην τρίτη θέση (την οποία δεν έχουμε καλύψει ακόμα), οπότε η λέξη μάλλον είναι \_ α κ ο σ. Χρειάζομαι ένα σύμφωνο για την πρώτη θέση που δεν είναι "ε, ρ, μ, υ, λ, χ, π". Μια πιθανή λέξη είναι "θάκος" (το "θ" δεν έχει δοκιμαστεί).} \\
\texttt{εικασία: θάκος} \\
            }
        }
    }
    & & \\ \\

    \theutterance \stepcounter{utterance}  
    & & & \multicolumn{2}{p{0.3\linewidth}}{
        \cellcolor[rgb]{0.9,0.9,0.9}{
            \makecell[{{p{\linewidth}}}]{
                \texttt{\tiny{[GM$|$GM]}}
                \texttt{Error: The guessed word is not a valid word for this game.} \\
            }
        }
    }
    & & \\ \\

    \theutterance \stepcounter{utterance}  
    & & & \multicolumn{4}{p{0.6\linewidth}}{
        \cellcolor[rgb]{0.9,0.9,0.9}{
            \makecell[{{p{\linewidth}}}]{
                \texttt{\tiny{[P1$\langle$GM]}}
                \texttt{Η εικασία σας δεν είναι έγκυρη λέξη για αυτό το παιχνίδι. Παρακαλώ δοκιμάστε ξανά.} \\
\\ 
\texttt{Δώστε την απάντησή σας μόνο σε αυτή τη μορφή.} \\
\texttt{εξήγηση: λεπτομέρειες} \\
\texttt{εικασία: λέξη} \\
            }
        }
    }
     \\ \\

    \theutterance \stepcounter{utterance}  
    & \multicolumn{4}{p{0.6\linewidth}}{
        \cellcolor[rgb]{0.561,0.737,0.561}{
            \makecell[{{p{\linewidth}}}]{
                \texttt{\tiny{[P1$\rangle$GM]}}
                \texttt{εξήγηση: Η λέξη είναι \_ α \_ ο σ, με το "κ" κίτρινο (άρα μπαίνει στην τρίτη θέση, αφού είναι η μόνη κενή). Αποκλεισμένα γράμματα: ε, ρ, μ, υ, λ, χ, π. Δοκιμάζω ένα νέο σύμφωνο για την πρώτη θέση που δεν έχω χρησιμοποιήσει, όπως το "δ".} \\
\texttt{εικασία: δάκος} \\
            }
        }
    }
    & & \\ \\

    \theutterance \stepcounter{utterance}  
    & & & \multicolumn{2}{p{0.3\linewidth}}{
        \cellcolor[rgb]{0.9,0.9,0.9}{
            \makecell[{{p{\linewidth}}}]{
                \texttt{\tiny{[GM$|$GM]}}
                \texttt{Error: The guessed word is not a valid word for this game.} \\
            }
        }
    }
    & & \\ \\

    \theutterance \stepcounter{utterance}  
    & & & \multicolumn{4}{p{0.6\linewidth}}{
        \cellcolor[rgb]{0.9,0.9,0.9}{
            \makecell[{{p{\linewidth}}}]{
                \texttt{\tiny{[P1$\langle$GM]}}
                \texttt{Η εικασία σας δεν είναι έγκυρη λέξη για αυτό το παιχνίδι. Παρακαλώ δοκιμάστε ξανά.} \\
\\ 
\texttt{Δώστε την απάντησή σας μόνο σε αυτή τη μορφή.} \\
\texttt{εξήγηση: λεπτομέρειες} \\
\texttt{εικασία: λέξη} \\
            }
        }
    }
     \\ \\

    \theutterance \stepcounter{utterance}  
    & \multicolumn{4}{p{0.6\linewidth}}{
        \cellcolor[rgb]{0.561,0.737,0.561}{
            \makecell[{{p{\linewidth}}}]{
                \texttt{\tiny{[P1$\rangle$GM]}}
                \texttt{εξήγηση: Η λέξη έχει μορφή \_ α \_ ο σ. Το "κ" είναι κίτρινο και η μόνη διαθέσιμη θέση είναι η τρίτη, άρα \_ α κ ο σ. Αποκλεισμένα: ε, ρ, μ, υ, λ, χ, π, δ. Δοκιμάζω το "σάκος", παρόλο που το πρώτο "σ" είναι διαφορετικής θέσης από το τελευταίο (η λέξη μπορεί να έχει ίδιο γράμμα σε διαφορετικές θέσεις).} \\
\texttt{εικασία: σάκος} \\
            }
        }
    }
    & & \\ \\

    \theutterance \stepcounter{utterance}  
    & & & \multicolumn{2}{p{0.3\linewidth}}{
        \cellcolor[rgb]{0.9,0.9,0.9}{
            \makecell[{{p{\linewidth}}}]{
                \texttt{\tiny{[GM$|$GM]}}
                \texttt{attempts = 5} \\
\texttt{target\_word = σακοσ} \\
\texttt{guess = σακοσ} \\
\texttt{guess\_feedback = σ<green> α<green> κ<green> ο<green> σ<green>} \\
            }
        }
    }
    & & \\ \\

    \theutterance \stepcounter{utterance}  
    & & & \multicolumn{2}{p{0.3\linewidth}}{
        \cellcolor[rgb]{0.9,0.9,0.9}{
            \makecell[{{p{\linewidth}}}]{
                \texttt{\tiny{[GM$|$GM]}}
                \texttt{game\_result = WIN} \\
            }
        }
    }
    & & \\ \\

\end{supertabular}
}

\end{document}
